\chapter{Regressão Quantílica}
\label{cap:qreg}
% ── índice ──────────────────────────────────────────────────────
\index{regressão quantílica|textbf}
\index{qreg@\texttt{qreg()}|textbf}
\index{mediana}
\index{quantil}
\index{heterocedasticidade}

% ─────────────────────────────────────────────────────────────────────────────
\section{O modelo}

O OLS estima $E[y \mid x]$ — a média condicional. A Regressão Quantílica
(QR) estima $Q_\tau[y \mid x]$ — o $\tau$-ésimo quantil condicional para
qualquer $\tau \in (0,1)$:

\[
  Q_\tau(y \mid x) = x'\beta(\tau)
\]

O vetor $\beta(\tau)$ varia com $\tau$: cada quantil tem seu próprio conjunto
de coeficientes. Com um único modelo é possível descrever a distribuição
condicional inteira de $y$.

% ─────────────────────────────────────────────────────────────────────────────
\section{O estimador}

O estimador minimiza a \textbf{função de perda assimétrica} (check function):

\[
  \hat\beta(\tau) = \arg\min_\beta \sum_{i=1}^{n} \rho_\tau(y_i - x_i'\beta)
\]

onde
\[
  \rho_\tau(u) =
  \begin{cases}
    \tau\, u       & \text{se } u \geq 0 \\
    (\tau - 1)\,u  & \text{se } u < 0
  \end{cases}
\]

Para $\tau = 0{,}5$: $\rho_{0,5}(u) = |u|/2$ — minimização da soma dos
desvios absolutos (LAD/mediana). Para $\tau \neq 0{,}5$, resíduos positivos
e negativos recebem pesos assimétricos, deslocando a estimativa para o
$\tau$-ésimo quantil.

\begin{nota}
  Ao contrário do OLS, a QR não tem solução de forma fechada — o problema
  é um programa linear resolvido por simplex ou métodos de ponto interior.
  Por isso a inferência padrão usa bootstrap (\lstinline{boot=}).
\end{nota}

% ─────────────────────────────────────────────────────────────────────────────
\section{Vantagens sobre o OLS}

\begin{itemize}
  \item \textbf{Heterocedasticidade:} quando $\mathrm{Var}[y \mid x]$ varia
        com $x$, os slopes $\beta_j(\tau)$ diferem entre quantis. O OLS
        captura apenas a média ponderada implícita.

  \item \textbf{Robustez:} a função check é linear nos resíduos, não
        quadrática. Outliers em $y$ influenciam menos $\hat\beta(\tau)$
        do que $\hat\beta_{\mathrm{OLS}}$.

  \item \textbf{Distribuição assimétrica:} salários, renda e preços de
        ativos têm caudas pesadas à direita. A QR descreve efeitos nas
        caudas sem suposições sobre a forma da distribuição.
\end{itemize}

% ─────────────────────────────────────────────────────────────────────────────
\section{Verificação por DGP}

O DGP abaixo introduz heterocedasticidade multiplicativa: a dispersão
de $y$ cresce com $x$, de modo que o efeito de $x$ é maior nos quantis
superiores do que nos inferiores.

\[
  y_i = 2{,}0 + 1{,}5\,x_i + (1{,}0 + 0{,}8\,x_i)\,\varepsilon_i,
  \quad \varepsilon_i \sim \mathcal{N}(0,1)
\]

O quantil condicional $Q_\tau(y \mid x)$ é crescente em $x$, e a
inclinação $\partial Q_\tau / \partial x$ aumenta com $\tau$: indivíduos
nos quantis superiores são mais sensíveis a $x$ do que os do quantil
inferior.

\begin{lstlisting}[caption={DGP heterocedástico — OLS vs.\ quatro quantis}]
seed(42)
input df
id
1
end
for i in 1..=10 { df = append(df, df) }
df |> filter(_n <= 1000)
generate df x = rnormal()
generate df e = rnormal()
generate df y = 2.0 + 1.5*x + (1.0 + 0.8*x)*e
let m_ols = ols(y ~ x, df)
let m_q25 = qreg(y ~ x, df, tau=0.25)
let m_q50 = qreg(y ~ x, df, tau=0.50)
let m_q75 = qreg(y ~ x, df, tau=0.75)
esttab(m_ols, m_q25, m_q50, m_q75)
\end{lstlisting}

\begin{lstlisting}[language={}, basicstyle=\ttfamily\small,
  frame=none, numbers=none, backgroundcolor=\color{codbg},
  xleftmargin=0pt, xrightmargin=0pt, breaklines=false]
                            (1)              (2)              (3)              (4)
                            OLS     QReg(τ=0.25)     QReg(τ=0.50)     QReg(τ=0.75)
──────────────────────────────────────────────────────────────────────────────────
x                     1.9947***        1.8739***        2.0149***        2.2462***
                       (0.0454)         (0.0687)         (0.0833)         (0.0579)
const                                  2.1659***        2.4525***        2.6637***
                                        (0.0302)         (0.0433)         (0.0369)
Constant              2.4538***
                       (0.0261)
──────────────────────────────────────────────────────────────────────────────────
N                          1000                0                0                0
R²                       0.6594
Adj. R²                  0.6591
──────────────────────────────────────────────────────────────────────────────────
* p<0.10  ** p<0.05  *** p<0.01
\end{lstlisting}

O slope de $x$ cresce monotonicamente: $1{,}87$ no Q25, $2{,}01$ no Q50
e $2{,}25$ no Q75. O OLS ($1{,}99$) é uma média ponderada implícita desses
efeitos heterogêneos — resume em um único número o que a QR decompõe em
toda a distribuição condicional.

A constante também aumenta com $\tau$ ($2{,}17 \to 2{,}45 \to 2{,}66$),
refletindo que os quantis superiores de $y$ estão acima da média para
qualquer valor de $x$.

% ─────────────────────────────────────────────────────────────────────────────
\section{Sintaxe}

\begin{lstlisting}
qreg(y ~ x1 + x2, df, tau=0.5)
qreg(y ~ x1 + x2, df, tau=0.5, boot=500)
\end{lstlisting}

O parâmetro \lstinline{tau} define o quantil desejado ($\tau \in (0,1)$,
padrão $0{,}5$). O parâmetro \lstinline{boot} ativa bootstrap para os
erros padrão (padrão 200 réplicas).

\begin{lstlisting}[caption={Fluxo típico — salários em múltiplos quantis}]
load "pnad.csv" as df

let q10 = qreg(sal ~ edu + exp, df, tau=0.10, boot=500)
let q50 = qreg(sal ~ edu + exp, df, tau=0.50, boot=500)
let q90 = qreg(sal ~ edu + exp, df, tau=0.90, boot=500)
let ols = ols(sal ~ edu + exp, df)

esttab(ols, q10, q50, q90)
\end{lstlisting}

\section{Capturando o resultado}

\begin{lstlisting}
let m = qreg(y ~ x, df, tau=0.75)
display m
let yhat = predict(m, "xb")
\end{lstlisting}
